\documentclass[a4paper,12pt]{article}
\usepackage[brazilian]{babel}
\usepackage[utf8]{inputenc}
\usepackage[T1]{fontenc}
\usepackage{epsfig}
\usepackage{array}

\usepackage{graphicx, url}
\usepackage[usenames]{color}
\usepackage{rotating}
\usepackage{url}
\usepackage{subfigure}
\usepackage{xspace}
\usepackage[verbose, dvips=true, pdftex=false, vtex=false, paperwidth=210mm,paperheight=295mm,top=35mm,left=30mm,bottom=35mm,right=30mm]{geometry}
\usepackage[authoryear]{natbib}
\usepackage{fancyhdr}
\usepackage{setspace}


\DeclareGraphicsExtensions{.png,.jpg,.pdf,.mps}
\bibliographystyle{icmc2}

\clubpenalty=10000
\widowpenalty=10000
\newcommand{\rev}[1]{{\color{red}#1}}

\newcommand{\titlehd}{\sffamily\bfseries\selectfont \scriptsize Uma Abordagem de Reestruturação de Sistemas Baseada em Requisitos de Qualidade Pré-Estabelecidos }

\pagestyle{fancy}
\headheight 61pt
\fancyhead{}
\fancyhead[LO]{\epsfig{file=icmc.jpg,height=15mm}}
\fancyhead[RO]{}
\fancyfoot{}

\begin{document}

\begin{center}
\vspace*{1cm}
{\fontsize{17.28}{20}\sffamily\bfseries\selectfont Uma Abordagem de Reestruturação de Sistemas Baseada em Requisitos de Qualidade Pré-Estabelecidos \\\ \\

Proposta de Pesquisa para Doutorado\\ \ \\
\vspace{3cm}}


\begin{table}[!th]
  \begin{center}
    \begin{tabular}{ll}

	  Orientador:      &  Prof. Dr. Márcio Eduardo Delamaro     \\
      Endereço:        &  Av. do Trabalhador São-Carlense, 400                    \\
	                   &  13560-970 - São Carlos/SP                               \\
      Email:           &  delamaro@icmc.usp.br                                     \\
	  
	  \vspace{0.6cm}
      Telefone:        &  (16) 3373-9669                                           \\

	  Candidato:       & Rafael Serapilha Durelli                             \\
	  Endereço:        & Rua Jesuino de Arruda, 1365, AP  403, Bloco 8 \\
                       & 13564-380 - São Carlos/SP                                 \\
      Email:           & rafa.durelli@gmail.com                                      \\
      Telefone:        & (16) 3415-6612                                            \\ \\
      \vspace{1cm}

    \end{tabular}
  \end{center}
\end{table}

\vfill
{\fontsize{18}{20}\sffamily\bfseries\selectfont Janeiro de 2012}

\end{center}

\vspace{3cm}

\pagebreak
\fancyhead[RO]{\titlehd}
\fancyfoot{\centerline{\thepage}}
\setcounter{page}{1}

\parskip 6pt

\doublespacing
\begin{abstract}
 Relatório Científico Parcial apresentado à FAPESP (Fundação de Amparo à Pesquisa do Estado de São Paulo) com o objetivo de elucidar as atividades realizadas  
 pelo bolsista Rafael Serapilha Durelli durante o primeiro período de vigência da bolsa concebida sob o Processo Número 2012/05168-4. 
 O referido período teve inicio em Junho de 2012 e foi finalizado em Março de 2013. 
 Além disso, este relatório também descreve as atividades que estão em andamentos, bem como as atividades a serem realizadas no próximo período. 
 Vale ressaltar que a partir de resultados preliminares alguns artigos científicos foram elaborados e publicados.
%Software é um produto que tende a evoluir constantemente com o intuito de satisfazer às necessidades de seus usuários. Para isso é necessário submetê-lo a constantes atividades de manutenção, que podem ocasionar um possível degradação no código-fonte. Dessa forma, o objetivo desta proposta de doutorado é o desenvolvimento de uma abordagem que auxilie o engenheiro de software a escolher o melhor conjunto de transformações para um sistema legado com base em requisitos de qualidade pré-estabelecidos e com isso possa identificar o impacto que tais transformações irão ocasionar no sistema legado. Isto é, a abordagem deverá indicar conjuntos de possíveis transformações para a evolução do sistema legado, além disso, deverá informar ao engenheiro o quanto tais transformações irão afetar a evolução/estrutura desse sistema. Com isso, o engenheiro de software poderia previamente escolher qual o melhor conjunto de  transformação se aplica dado um requisito de qualidade pré-estabelecido. Após a escolha das transformações a abordagem também deverá realizar automaticamente tais transformações a fim de reestruturar o sistema legado em um novo sistema. Será avaliado a abordagem por meio de experimentos controlados.

\end{abstract}


      



\section{Introdução e Contexto}


\subsection{Motivação}\label{mot}


\subsection{Objetivos}\label{obj}


\subsection{Organização da Proposta}\label{org}

       

\section{Fundamentação Teórica}\label{sec:fundamentacao}


\subsection{Mineração de Padrões de Projeto}\label{sub:mineracao}


 
\subsection{Mineração de Interesses Transversais}\label{sub:mineracao}


\section{Proposta}\label{sec:proposta}



\subsection{Objetivos}\label{sec:Objetivos}


\subsection{Desafios de Pesquisa com Relação ao Projeto}\label{desafios}

\subsection{Metodologia}\label{sec:metodologia}

   
\subsubsection{Atividades e Cronograma}\label{sec:atividadeECronograma}


\subsection{Avaliação}\label{sec:avaliacao}

\subsection{Resultados Esperados}\label{sec:ResultadosEsperados}

\section{Trabalhos Relacionados}\label{sec:relacionados}
 
     

\singlespacing 
\bibliographystyle{icmc}
\bibliography{referencias2}

\end{document}
